\documentclass[12pt,a4paper]{article}
\usepackage[utf8]{inputenc}
\usepackage{amsmath}
\usepackage{amssymb}
\usepackage{amsthm}
\usepackage{geometry}

% Page setup
\geometry{margin=1in}

\title{\textbf{Rigorous Derivation of DNF and Tseitin Encoding}}
\author{}
\date{}

\begin{document}

\maketitle

\section*{Problem Statement}

Consider the formula:
\[
    \varphi \equiv x_1 \lor \neg(x_2 \land \neg(x_3 \leftrightarrow (x_1 \to \neg x_2)))
\]

\textbf{(a)} Give a DNF formula that is semantically equivalent to $\varphi$, using techniques studied in class.

\textbf{(b)} Use Tseitin encoding to write an equisatisfiable formula $\psi$ in conjunctive normal form. Would $\neg \psi$ be also equisatisfiable with $\neg \varphi$? Give justification for your answer.

\section*{Part (a): Disjunctive Normal Form (DNF)}

\textbf{Solution:}

We derive the Disjunctive Normal Form (DNF) by systematically applying logical equivalences.

\subsection*{Step 1: Eliminate Implications and Equivalences}
First, we address the innermost implication $x_1 \to \neg x_2$. Using the equivalence $A \to B \equiv \neg A \lor B$:
\[ (x_1 \to \neg x_2) \equiv (\neg x_1 \lor \neg x_2) \]
Substituting this back into $\varphi$:
\[ \varphi \equiv x_1 \lor \neg(x_2 \land \neg(x_3 \leftrightarrow (\neg x_1 \lor \neg x_2))) \]

Next, we address the equivalence $A \leftrightarrow B \equiv (A \land B) \lor (\neg A \land \neg B)$. Let $Q = (\neg x_1 \lor \neg x_2)$.
\[ x_3 \leftrightarrow Q \equiv (x_3 \land Q) \lor (\neg x_3 \land \neg Q) \]
Using De Morgan's Law, $\neg Q = \neg(\neg x_1 \lor \neg x_2) = (x_1 \land x_2)$. Thus the equivalence $E$ becomes:
\[ E \equiv (x_3 \land (\neg x_1 \lor \neg x_2)) \lor (\neg x_3 \land (x_1 \land x_2)) \]

\subsection*{Step 2: Push Negations Inward}
The formula is now $\varphi \equiv x_1 \lor \neg(x_2 \land \neg E)$. By De Morgan's Law:
\[ \varphi \equiv x_1 \lor (\neg x_2 \lor \neg \neg E) \]
By Double Negation Elimination:
\[ \varphi \equiv x_1 \lor \neg x_2 \lor E \]
Substituting $E$ back in:
\[ \varphi \equiv x_1 \lor \neg x_2 \lor \left[ (x_3 \land (\neg x_1 \lor \neg x_2)) \lor (\neg x_3 \land x_1 \land x_2) \right] \]

\subsection*{Step 3: Distribute to obtain DNF}
We apply the distributive law to $(x_3 \land (\neg x_1 \lor \neg x_2))$:
\[ (x_3 \land \neg x_1) \lor (x_3 \land \neg x_2) \]
Substituting this back:
\[ \varphi \equiv x_1 \lor \neg x_2 \lor (x_3 \land \neg x_1) \lor (x_3 \land \neg x_2) \lor (\neg x_3 \land x_1 \land x_2) \]

\subsection*{Step 4: Simplification}
Using the Absorption Law $A \lor (A \land B) \equiv A$:

\noindent \textbf{Term 1:} $x_1$ absorbs $(\neg x_3 \land x_1 \land x_2)$.

\noindent \textbf{Term 2:} $\neg x_2$ absorbs $(x_3 \land \neg x_2)$.

\noindent Reduced DNF:
\[ \varphi \equiv x_1 \lor \neg x_2 \lor (x_3 \land \neg x_1) \]
Using $A \lor (\neg A \land B) \equiv A \lor B$:
\[ x_1 \lor (\neg x_1 \land x_3) \equiv x_1 \lor x_3 \]
Final DNF:
\[ \varphi_{\text{DNF}} \equiv x_1 \lor \neg x_2 \lor x_3 \]

\newpage

\section*{Part (b): Tseitin Encoding}

\textbf{Solution:}

\subsection*{1. Tseitin Construction}
We introduce fresh variables $y_i$ for sub-expressions.
\begin{align*}
    y_1 &\leftrightarrow \neg x_2 \\
    y_2 &\leftrightarrow (x_1 \to y_1) \\
    y_3 &\leftrightarrow (x_3 \leftrightarrow y_2) \\
    y_4 &\leftrightarrow \neg y_3 \\
    y_5 &\leftrightarrow (x_2 \land y_4) \\
    y_6 &\leftrightarrow \neg y_5 \\
    y_7 &\leftrightarrow (x_1 \lor y_6)
\end{align*}

\subsection*{2. CNF Clauses for Definitions}

\noindent \textbf{1. Definition of $y_1$:} $(y_1 \leftrightarrow \neg x_2)$
\[ (\neg y_1 \lor \neg x_2) \land (x_2 \lor y_1) \]

\noindent \textbf{2. Definition of $y_2$:} $(y_2 \leftrightarrow (\neg x_1 \lor y_1))$
\[ (\neg y_2 \lor \neg x_1 \lor y_1) \land (x_1 \lor y_2) \land (\neg y_1 \lor y_2) \]

\noindent \textbf{3. Definition of $y_3$:} $(y_3 \leftrightarrow (x_3 \leftrightarrow y_2))$
\[
\begin{aligned}
(\neg y_3 \lor \neg x_3 \lor y_2) &\land (\neg y_3 \lor x_3 \lor \neg y_2) \\
&\land (y_3 \lor x_3 \lor y_2) \land (y_3 \lor \neg x_3 \lor \neg y_2)
\end{aligned}
\]

\noindent \textbf{4. Definition of $y_4$:} $(y_4 \leftrightarrow \neg y_3)$
\[ (\neg y_4 \lor \neg y_3) \land (y_3 \lor y_4) \]

\noindent \textbf{5. Definition of $y_5$:} $(y_5 \leftrightarrow (x_2 \land y_4))$
\[ (\neg y_5 \lor x_2) \land (\neg y_5 \lor y_4) \land (\neg x_2 \lor \neg y_4 \lor y_5) \]

\noindent \textbf{6. Definition of $y_6$:} $(y_6 \leftrightarrow \neg y_5)$
\[ (\neg y_6 \lor \neg y_5) \land (y_5 \lor y_6) \]

\noindent \textbf{7. Definition of $y_7$:} $(y_7 \leftrightarrow (x_1 \lor y_6))$
\[ (\neg y_7 \lor x_1 \lor y_6) \land (\neg x_1 \lor y_7) \land (\neg y_6 \lor y_7) \]

\noindent \textbf{8. Root Unit Clause:}
\[ (y_7) \]

\subsection*{3. Equisatisfiability of Negations}
\textbf{Question:} Is $\neg \psi$ equisatisfiable with $\neg \varphi$? \\
\textbf{Answer:} No.

\textbf{Justification:}
Two formulas $A$ and $B$ are equisatisfiable if ($A$ is satisfiable $\iff$ $B$ is satisfiable).

\noindent \textbf{Case 1:} $\neg \varphi$ is satisfiable only if $\varphi$ is \textbf{not} a tautology.

\noindent \textbf{Case 2:} $\neg \psi$ is equivalent to $\neg y_{\text{root}} \lor \neg (\text{Definitions})$. Since it is always possible to violate the definitions (by choosing variable assignments that make the gates False), $\neg \psi$ is \textbf{always satisfiable}.

Because $\neg \psi$ is always satisfiable regardless of the logical validity of $\varphi$, they are not equisatisfiable.

\end{document}